\subsection{Routing-time estimates}
\label{app:routing-proof}
Write $q_1=\log(N-1)$; $N\ge4$ implies $q_1>1$.
The quantitative bounds in parts (i)--(ii) of \cref{thm:separation} are
\begin{align*}
 \frac{N-1}{2e\gamma}\log\!\left(1+\frac{2}{\varepsilon^2}\right)
 &\le\Tr(\Dlaw1)\le\frac{4(N-1)}{\gamma}\left[\log(1+\varepsilon^{-2})+\frac{1}{e}\right],\\
 \frac{N-1}{5\gamma \varepsilon}
 &\le\Tr(\Dlaw0)\le\frac{4(N-1)}{\gamma}\left[\frac{\sqrt2\pi}{\varepsilon}+\frac{8}{e}\right].
\end{align*}

\begin{lemma}[Attention bounds on the routing interval]
\label{lem:attention-bounds}
For $q\in[0,q_1]$,
\begin{equation}
 \frac{e^q}{2(N-1)}\le \alpha(q)\le\frac{e^q}{N},\qquad
 \frac{1}{N}\le \alpha(q)\le\frac{1}{2},\qquad
 \phi'(q)\ge\frac{e^q}{4(N-1)}.
 \label{eq:alpha-bounds}
\end{equation}
For $\kappa\in\{0,1\}$ and $q\in[0,1]$, $\phi'(q)\le \tfrac{e}{N-1}$.
\end{lemma}
\begin{proof}
\begin{align*}
 1\le e^q\le N-1
 &\quad\Longrightarrow\quad
 N\le e^q+N-1\le2(N-1),\\
 \phi'(q)=(1+\rho)\alpha(q)(1-\alpha(q))
 &\ge\frac{\alpha(q)}{2}\ge\frac{e^q}{4(N-1)}&&\text{($\rho\ge0$, $\alpha\le\tfrac12$)},\\
 \kappa\in\{0,1\},\ 0\le q\le1:\quad
 \phi'(q)&\le(1+\rho)\alpha(q)
 \le\frac{N}{N-1}\frac{e^q}{N}\le\frac{e}{N-1}.\qedhere
\end{align*}
\end{proof}

\begin{lemma}[Primitive bounds]
\label{lem:primitive-bounds}
For $0\le q\le q_1$, the two laws satisfy
\begin{align}
 \kappa=1:\quad &q\le F(q)\le2q,\label{eq:ordinary-primitive-bounds}\\
 \kappa=0:\quad &\frac{q^2}{8}\le F(q)\le2q^2\quad(0\le q\le1),\qquad
 F(q)\ge\frac{1}{8}\quad(1\le q\le q_1).
 \label{eq:cancelled-primitive-bounds}
\end{align}
\end{lemma}
\begin{proof}
\proofstep{(a) Ordinary law.}

Here $\rho=0$, $\phi=\alpha$, and $\phi'(q)=\alpha(1-\alpha)$.
\Cref{lem:attention-bounds} gives $1\le \tfrac{\phi}{\phi'(q)}\le2$;
$F(q)=\int_0^q\tfrac{\phi(v)}{\phi'(v)}\,dv\in[q,2q]$.

\proofstep{(b) Cancelled law.}

Here $\rho=\tfrac{1}{N-1}$, and
\[
 \dv{F(q)}{q}=\frac{\alpha(q)-\frac{1}{N}}{\alpha(q)(1-\alpha(q))}\ge0\quad(q\ge0).
\]
For $0\le q\le1$, \cref{lem:attention-bounds} and convexity give
\begin{align*}
 \alpha(q)-\frac{1}{N}
 &=\frac{(e^q-1)(N-1)}{N(e^q+N-1)}\le\frac{e^q-1}{N},\\
 \alpha(q)(1-\alpha(q))&\ge\frac{1}{2N},&
 e^q-1&\le(e-1)q\le2q,\\
 \dv{F(q)}{q}&\le4q,&F(q)&\le2q^2.
\end{align*}
For $0\le q\le q_1$,
\begin{align*}
 \dv{F(q)}{q}&\ge\frac{\alpha(q)-\frac{1}{N}}{\alpha(q)}
 =\frac{N-1}{N}(1-e^{-q})\ge\frac{1}{2}(1-e^{-q}),\\
 0\le q\le1:\quad
 1-e^{-q}&\ge q-\frac{q^2}{2}\ge \frac{q}{2},\\
 F(q)&\ge\int_0^q \frac{v}{4}\,dv=\frac{q^2}{8},\\
 q\ge1:\quad F(q)&\ge F(1)\ge\frac{1}{8}&&\text{($F$ is nondecreasing)}.\qedhere
\end{align*}
\end{proof}

\begin{proof}[Proof of \cref{thm:separation}(i) and \cref{rem:rescaling}]
For a fixed displayed-prefix law and independent error sign, write $\mathcal L_\eta$ for correlation loss. Then
\[
 \mathcal L_\eta(\bm\theta)=(1-2\eta)\mathcal L_0(\bm\theta),\qquad
 \bm\theta_\eta(t)=\bm\theta_0((1-2\eta)t),\qquad
 T(\eta,\kappa)=\frac{T(0,\kappa)}{1-2\eta}.
\]
The trajectory identity follows from the chain rule and uniqueness in the original Euclidean coordinate metric.
It holds on corresponding gate-valid intervals, with infinite routing times interpreted as infinite.
Thus independent corruption at fixed $\kappa=1$ strictly increases finite routing time for $\eta>0$;
the failure threshold at fixed $\kappa<0$ is unchanged.

By \cref{lem:primitive-bounds}, $\varepsilon^2+F(q)>0$ on $[0,q_1]$.
\Cref{lem:invariant,lem:attention-bounds,lem:primitive-bounds} give
\begin{align*}
 \Tr
 &\le\frac{4(N-1)}{\gamma}\int_0^{q_1}\frac{e^{-q}}{\varepsilon^2+q}\,dq\\
 &\le\frac{4(N-1)}{\gamma}
 \left(\int_0^1\frac{dq}{\varepsilon^2+q}+\int_1^\infty e^{-q}\,dq\right)\\
 &=\frac{4(N-1)}{\gamma}\left(\log(1+\varepsilon^{-2})+e^{-1}\right),\\[2pt]
 \Tr
 &\ge\frac{N-1}{e\gamma}\int_0^1\frac{dq}{\varepsilon^2+2q}
 =\frac{N-1}{2e\gamma}\log(1+2\varepsilon^{-2}).\qedhere
\end{align*}
\end{proof}

\begin{proof}[Proof of \cref{thm:separation}(ii)]
By \cref{lem:primitive-bounds}, $\varepsilon^2+F(q)\ge \varepsilon^2>0$.
\Cref{lem:invariant,lem:attention-bounds,lem:primitive-bounds} give
\begin{align*}
 \Tr
 &\ge\frac{N-1}{e\gamma}\int_0^1\frac{dq}{\varepsilon^2+2q^2}
 =\frac{N-1}{e\gamma}\frac{\arctan(\frac{\sqrt2}{\varepsilon})}{\sqrt2\varepsilon}
 \ge\frac{N-1}{\gamma \varepsilon}\frac{\arctan\sqrt2}{\sqrt2e},\\
 \arctan\sqrt2
 &=\int_0^1\frac{dv}{1+v^2}+\int_1^{\sqrt2}\frac{dv}{1+v^2}\\
 &\ge\int_0^1(1-v^2)\,dv+\frac{\sqrt2-1}{3}=\frac{1+\sqrt2}{3},\\
 \sqrt2<\frac{3}{2},\ e<\frac{11}{4}
 &\quad\Longrightarrow\quad
 \frac{\arctan\sqrt2}{\sqrt2e}>\frac{20}{99}>\frac{1}{5},\\[2pt]
 \Tr
 &\le\frac{4(N-1)}{\gamma}
 \left(\int_0^1\frac{dq}{\varepsilon^2+\frac{q^2}{8}}+8\int_1^{q_1}e^{-q}\,dq\right)\\
 &\le\frac{4(N-1)}{\gamma}
 \left(\int_0^\infty\frac{dq}{\varepsilon^2+\frac{q^2}{8}}+8e^{-1}\right)\\
 &=\frac{4(N-1)}{\gamma}
 \left(\frac{\sqrt2\pi}{\varepsilon}+\frac{8}{e}\right).\qedhere
\end{align*}
\end{proof}

\subsection{Diluted cancellation: proof of \texorpdfstring{\cref{cor:dilution}}{the dilution corollary}}
\label{app:dilution}
\begin{proof}[Proof of \cref{cor:dilution}]
For $0\le\lambda\le1$, \cref{def:Dkappa} draws $M\sim\mathrm{Bernoulli}(1-\lambda)$ anti-copies, so $\Dlaw\lambda=\lambda\Dlaw1+(1-\lambda)\Dlaw0$.
Write $\kappa=\lambda$, $\rho=\tfrac{1-\lambda}{N-1}$ and $\alpha_0=\tfrac{\rho}{1+\rho}=\tfrac{1-\lambda}{N-\lambda}$.
\proofstep{(a) The primitive.}
From \cref{eq:primitive}, $F'(q)=\tfrac{\alpha(q)-\alpha_0}{\alpha(q)(1-\alpha(q))}=\tfrac{g(q)}{1-\alpha(q)}$ with
\[
 g(q)=1-\frac{\alpha_0}{\alpha(q)}=\frac{N-1}{N-\lambda}\Bigl[\lambda+(1-\lambda)(1-e^{-q})\Bigr]\ge0 .
\]
On $[0,q_1]$, $1\le\tfrac{1}{1-\alpha(q)}\le2$, and $\tfrac{N-1}{N-\lambda}\in[\tfrac{3}{4},1]$. For $0\le q\le1$, $\tfrac{q}{2}\le1-e^{-q}\le q$, and
$\lambda+(1-\lambda)q\ge\tfrac{\lambda+q}{4}$ (if $\lambda\le\tfrac12$ the left side is at least $\lambda+\tfrac{q}{2}$; otherwise it is at least $\lambda>\tfrac12\ge\tfrac{\lambda+q}{4}$). Hence
\[
 \frac{3}{32}(\lambda+q)\le F'(q)\le2(\lambda+q),\qquad
 \frac{3}{32}\Bigl(\lambda q+\frac{q^2}{2}\Bigr)\le F(q)\le2\Bigl(\lambda q+\frac{q^2}{2}\Bigr)\qquad(0\le q\le1),
\]
and, since $F$ is nondecreasing, $F(q)\ge F(1)\ge\tfrac{3}{64}$ on $[1,q_1]$.
\proofstep{(b) Time bounds.}
\Cref{lem:attention-bounds} gives $\phi'(q)\ge\tfrac{e^q}{4(N-1)}$ on $[0,q_1]$, and its upper bound $\phi'(q)\le(1+\rho)\alpha(q)\le\tfrac{e}{N-1}$ on $[0,1]$ holds for every $\lambda\in[0,1]$ because $1+\rho\le\tfrac{N}{N-1}$.
Write $I=\int_0^1\tfrac{dq}{\varepsilon^2+\lambda q+q^2}$; note $I\ge\tfrac{1}{3}$. Using $\varepsilon^2+F\le2(\varepsilon^2+\lambda q+q^2)$ and $\varepsilon^2+F\ge\tfrac{3}{64}(\varepsilon^2+\lambda q+q^2)$ on $[0,1]$, \cref{eq:time} gives
\begin{align*}
 \Tr&\ge\frac{N-1}{e\gamma}\int_0^1\frac{dq}{\varepsilon^2+F(q)}\ge\frac{N-1}{2e\gamma}\,I,\\
 \Tr&\le\frac{4(N-1)}{\gamma}\left(\int_0^1\frac{e^{-q}\,dq}{\varepsilon^2+F(q)}+\int_1^{q_1}\frac{e^{-q}\,dq}{F(1)}\right)
 \le\frac{256(N-1)}{3\gamma}\bigl(I+e^{-1}\bigr)\le\frac{256(1+3e^{-1})(N-1)}{3\gamma}\,I .
\end{align*}
\proofstep{(c) The integral.}
Since $\max\{a,b,c\}\le a+b+c\le3\max\{a,b,c\}$, $I$ is within a factor $3$ of $J=\int_0^1\tfrac{dq}{\max\{\varepsilon^2,\lambda q,q^2\}}$.
If $\lambda\le\varepsilon$, then $\lambda q\le\max\{\varepsilon^2,q^2\}$ and $J=\tfrac{2}{\varepsilon}-1\in[\tfrac{1}{\varepsilon},\tfrac{2}{\varepsilon}]$.
If $\lambda>\varepsilon$, the three regimes meet at $q=\tfrac{\varepsilon^2}{\lambda}<\lambda\le1$, and
\[
 J=\frac{1}{\lambda}+\frac{2}{\lambda}\log\frac{\lambda}{\varepsilon}+\frac{1}{\lambda}-1
 \in\left[\frac{1+2\log(\lambda/\varepsilon)}{\lambda},\ \frac{2+2\log(\lambda/\varepsilon)}{\lambda}\right].
\]
With $C=\tfrac{1+\log(1+\lambda/\varepsilon)}{\lambda}$: for $\lambda\le\varepsilon$, $C\ge\tfrac{1}{\lambda}\ge\tfrac{1}{\varepsilon}$, so $\min\{\tfrac{1}{\varepsilon},C\}=\tfrac{1}{\varepsilon}$;
for $\lambda>\varepsilon$, $\log\tfrac{\lambda}{\varepsilon}\le\log(1+\tfrac{\lambda}{\varepsilon})\le\log\tfrac{\lambda}{\varepsilon}+\log2$ and $C\le\tfrac{1+\log2}{\varepsilon}$.
In both cases $J$, and hence $I$, is within absolute constants of $\min\{\tfrac{1}{\varepsilon},C\}$; with part (b) this proves the corollary.
\end{proof}
Numerically (\texttt{kappa\_sweep.py}), $F(q)/(\lambda q+\tfrac{q^2}{2})\in[0.67,1.14]$ on $[0,1]$ for $N\in\{8,16,64,256\}$, and
$\tfrac{\gamma\Tr}{(N-1)I}\in[1.03,1.51]$ over $N\in\{16,64\}$, $\varepsilon\in\{10^{-2},10^{-3},10^{-4}\}$ and eleven values of $\lambda\in[0,1]$, well inside the proved constants.

\subsection{Over-cancellation: exact barrier and critical equality}
\label{app:critical}

\begin{lemma}[Barrier location and size]
\label{lem:barrier}
For $-1\le \kappa<0$, let $\alpha_0=\tfrac{1+|\kappa|}{N+|\kappa|}$ and $q_0=\log(1-\kappa)$.
Then $\phi(q_0)=0$, $\alpha(q_0)=\alpha_0$, and
\begin{equation}
 \min_{q\ge0}F(q)=F(q_0)<0,\qquad
 \frac{9\kappa^2}{200}\le -F(q_0)\le2\kappa^2,\qquad
 \frac{1}{N}<\alpha_0\le\frac{2}{N}.
 \label{eq:barrier-bounds}
\end{equation}
\end{lemma}
\begin{proof}
\proofstep{(a) Location and sign.}
\begin{align*}
 \frac{(N-1)\alpha_0}{1-\alpha_0}&=(N-1)\rho=1-\kappa,\\
 \dv{F(q)}{q}&=\frac{\alpha(q)-\alpha_0}{\alpha(q)(1-\alpha(q))},\\
 q<q_0&\Rightarrow \dv{F(q)}{q}<0,&q>q_0&\Rightarrow \dv{F(q)}{q}>0,\\
 |\kappa|(N-2)\le N&\Rightarrow \alpha_0\le\frac{2}{N}.
\end{align*}
\proofstep{(b) Lower bound.}
\begin{align*}
 -F(q_0)
 &=\int_{\frac{1}{N}}^{\alpha_0}\frac{\alpha_0-\alpha}{\alpha^2(1-\alpha)^2}\,d\alpha
 \ge\frac{(\alpha_0-\frac{1}{N})^2}{2\alpha_0^2}
 =\frac{(N \alpha_0-1)^2}{2(N \alpha_0)^2},\\
 N \alpha_0-1&=\frac{|\kappa|(N-1)}{N+|\kappa|}\ge\frac{3}{5}|\kappa|,&N \alpha_0&\le2,\\
 -F(q_0)&\ge\frac{9\kappa^2}{200}.
\end{align*}
\proofstep{(c) Upper bound.}
For $\alpha\in[\tfrac{1}{N},\alpha_0]$,
\begin{align*}
 \alpha^2(1-\alpha)^2&\ge\frac{(1-\alpha_0)^2}{N^2}\ge\frac{1}{4N^2},\\
 -F(q_0)&\le4N^2\int_{\frac{1}{N}}^{\alpha_0}(\alpha_0-\alpha)\,d\alpha
 =2(N \alpha_0-1)^2\le2\kappa^2.\qedhere
\end{align*}
\end{proof}

The failure part also yields the following quantitative conclusions.
Below this barrier, $q(t)\uparrow q_z<q_0$ and $r(t)\to0$, where $q_z$ is the unique root of $\varepsilon^2+F(q_z)=0$.
At equality, $q(t)\uparrow q_0$ and $r(t)\downarrow0$.
Both failure cases have $\alpha(q(t))<\tfrac{2}{N}$ for every finite $t$.
Moreover $\tfrac{9\kappa^2}{200}\le-F(q_0)\le2\kappa^2$;
$\varepsilon\le\tfrac{|\kappa|}{9}$ suffices for failure, and $\varepsilon>\sqrt2|\kappa|$ suffices for routing.
\begin{proof}[Proof of \cref{thm:separation}(iii)]
\proofstep{(a) Above the barrier.}
If $\varepsilon^2>-F(q_0)$, then $\min_{[0,q_1]}(\varepsilon^2+F)>0$ by \cref{lem:barrier}; \cref{lem:invariant} gives $\Tr<\infty$.
\proofstep{(b) Below the barrier.}

\Cref{lem:barrier} gives strict decrease of $F$ on $[0,q_0]$, with
$\varepsilon^2+F(0)>0>\varepsilon^2+F(q_0)$; hence $q_z$ exists uniquely.
By \cref{lem:invariant},
\[
 r(t)>0,\qquad r(t)^2=\varepsilon^2+F(q(t)),\qquad\ddt{q}>0
 \quad\Longrightarrow\quad 0\le q(t)<q_z.
\]
Let $q_\infty=\lim_{t\to\infty}q(t)\le q_z$.
If $q_\infty<q_z$, then
\[
 \ddt{q}(t)\ge\min_{0\le q\le q_\infty}
 \gamma(\varepsilon^2+F(q))\phi'(q)>0,
\]
so $q(t)\to\infty$, a contradiction. Therefore $q_\infty=q_z$ and
$r(t)^2\to \varepsilon^2+F(q_z)=0$.

\proofstep{(c) Equality.}

At $q=q_0$, $\phi=0$ and
\begin{align*}
 \left.\dv{F(q)}{q}\right|_{q=q_0}&=0,\\
 \left.\frac{d^2F(q)}{dq^2}\right|_{q=q_0}
 &=\left.\dv{}{q}\left(\frac{\phi(q)}{\phi'(q)}\right)\right|_{q=q_0}\\
 &=\left.\left[1-\frac{\phi(q)\phi''(q)}{\bigl(\phi'(q)\bigr)^2}\right]\right|_{q=q_0}=1.
\end{align*}
Taylor's theorem gives $b>0$ such that, for $q_0-b\le q\le q_0$,
\[
 \frac{1}{4}(q-q_0)^2\le F(q)-F(q_0)\le\frac{3}{4}(q-q_0)^2.
\]
On $[0,q_0]$, $\phi'(q)\le\tfrac{1+\rho}{4}$, so
\[
 \int_{q_0-b}^{q_0}\frac{dq}{\gamma[\varepsilon^2+F(q)]\phi'(q)}
 \ge\frac{16}{3\gamma(1+\rho)}\int_{q_0-b}^{q_0}\frac{dq}{(q-q_0)^2}=\infty.
\]
Thus $q(t)<q_0$ at every finite time.
The same bounded-monotone limit argument as in part (b) yields
$q(t)\uparrow q_0$ and $r(t)^2\to \varepsilon^2+F(q_0)=0$.
Since $\phi(q)<0$ for $q<q_0$, $\ddt{r}<0$.
\proofstep{(d) Attention and sufficient conditions.}
Finally $q_0=\log(1-\kappa)\le\log2<\log(N-1)=q_1$, so routing fails. In both failure cases $\alpha(q(t))<\alpha(q_0)=\alpha_0\le\tfrac{2}{N}$.
The bounds in \cref{lem:barrier} give
\[
 \varepsilon\le\frac{|\kappa|}{9}\ \Longrightarrow\ \varepsilon^2<\frac{9\kappa^2}{200}\le-F(q_0),\qquad
 \varepsilon>\sqrt2|\kappa|\ \Longrightarrow\ \varepsilon^2>2\kappa^2\ge-F(q_0).\qedhere
\]
\end{proof}
